% ID: SEL-D0045-OB-E03
% ESTADO: cotejada. Segunda pasada visual y borrador compilado revisados.
% Fuente: Selectividad/sel_2014_matematicas_aplicadas/EXAMEN -ANDAL-5.pdf
% SHA-256: fe436fc3af8e56d0694316005aedeca2fffefbeb445aa9f60efa0aa6da4cba25
% Página: 2; opción B; ejercicio original 3.

El 25\% de los estudiantes de una Universidad lee las noticias en prensa escrita en papel, el 70\% en prensa digital y el 10\% en ambos formatos. Elegido, al azar, un estudiante de esa Universidad:
\begin{enumerate}[label=\alph*)]
\item \textbf{(1 punto)} Calcule la probabilidad de que lea las noticias en formato papel o digital.
\item \textbf{(0.75 puntos)} Sabiendo que lee las noticias en prensa digital, calcule la probabilidad de que también las lea en prensa escrita en papel.
\item \textbf{(0.75 puntos)} \textquestiondown{}Cuál es la probabilidad de que lea las noticias exclusivamente en uno de los dos formatos?
\end{enumerate}
